\section{Fourier Transform}

\subsection{Definition}

\begin{align*}
    \mathcal{F}\left\{f(x)\right\}           & = F(\alpha) = \frac{1}{2\pi} \int_{-\infty}^\infty f(x)e^{-i\alpha x}~dx \\
    \mathcal{F}^{-1}\left\{F(\alpha)\right\} & = f(x) = \int_{-\infty}^\infty F(\alpha)e^{i\alpha x}~d\alpha
\end{align*}

where $f(x) \to 0$ as $x \to \pm\infty$.

\vsp

Other definitions exist, such as that obtained by swapping the $1/(2\pi)$ factor between the two equations or using $1/\sqrt{2\pi}$ for both.

\vsp

The Fourier transform essentially takes the limit $L \to \infty$ on the Fourier series, lifting the constraint of a $2L$-periodic function. That is, the Fourier transform is a generalization of the Fourier series and can be used on non-periodic functions as well. This is useful for problems such as the wave equation with an infinitely long string. Furthermore, the Fourier transform is inherently continuous due to $\alpha$ whereas the Fourier series is discrete due to $n$.

\subsection{Properties}

\subsubsection{Differentiation}

The Fourier transform turns differentiation into multiplication:

\begin{align*}
    \mathcal{F}\left\{f'(x)\right\}      & = i\alpha F(\alpha) \tag{1st derivative}     \\
    \mathcal{F}\left\{f^{(n)}(x)\right\} & = (i\alpha)^n F(\alpha) \tag{nth derivative} \\
\end{align*}

\subsubsection{Integration}

The Fourier transform turns integration into division:

\begin{equation*}
    \mathcal{F}\left\{\int_{-\infty}^x f(s)~ds\right\} = \frac{F(\alpha)}{i\alpha} + \pi F(0)\delta(\alpha)
\end{equation*}

The $\pi F(0)\delta(\alpha)$ term represents a constant spike at $\alpha = 0$.

\subsection{Solving the Wave Equation for an Infinite String}

The Fourier transform can be used to solve the wave equation for an infinite string, which is given by

\begin{equation*}
    u_{tt} - c^2 u_{xx} = 0\text{, for } -\infty < x < \infty\text{, } t > 0
\end{equation*}

with initial conditions

\begin{equation*}
    u(x,0) = f(x)\text{, } u_t(x,0) = g(x)
\end{equation*}

where $u(x,t)$ is the displacement of the string at position $x$ and time $t$, and $c$ is the wave speed.

\vsp

Represent the solution $u(x,t)$ as an inverse Fourier transform:

\begin{equation*}
    u(x,t) = \int_{-\infty}^\infty U(\alpha,t)e^{i\alpha x}~d\alpha
\end{equation*}

Differentiate with respect to $t$ and $x$:

\begin{align*}
    u_{tt} & = \int_{-\infty}^\infty U''(\alpha,t)e^{i\alpha x}~d\alpha         \\
    u_{xx} & = -\int_{-\infty}^\infty \alpha^2 U(\alpha,t)e^{i\alpha x}~d\alpha
\end{align*}

Note that $U(\alpha, t)$ is strictly in terms of a single independent variable $t$ because $\alpha$ is just a dummy variable.

\vsp

Substitute these into the wave equation to yield a second-order ODE in terms of $t$:

\begin{align*}
    \int_{-\infty}^\infty \underbrace{\left[U''(\alpha,t) + \alpha^2 c^2 U(\alpha,t) \right]}_{=0} e^{i\alpha x}~d\alpha = 0 \\
    U''(\alpha,t) + \alpha^2 c^2 U(\alpha,t) = 0
\end{align*}

The initial conditions are equivalently

\begin{align*}
    U(\alpha,0)  & = F(\alpha) \\
    U'(\alpha,0) & = G(\alpha)
\end{align*}

\vsp

Thus, the general solution for $U(\alpha,t)$ is

\begin{align*}
    U(\alpha,t) = c_1(\alpha)\cos(c\alpha t) + c_2(\alpha)\sin(c\alpha t) \\
    U(\alpha,t) = A(\alpha)\cos(c\alpha t) + B(\alpha)\frac{\sin(c\alpha t)}{c\alpha}
\end{align*}

\vsp

where $A(\alpha) = c_1(\alpha)$ and $B(\alpha) = c\alpha~c_2(\alpha)$ are arbitrary functions of $\alpha$ that are determined by the initial conditions.

\vsp

Applying the initial conditions:

\begin{align*}
    U(\alpha,0) = F(\alpha) = A(\alpha) \\
    U'(\alpha,0) = G(\alpha) = B(\alpha)
\end{align*}

\vsp

Thus, the solution for $U(\alpha,t)$ is

\begin{equation*}
    U(\alpha,t) = F(\alpha)\cos(c\alpha t) + G(\alpha)\frac{\sin(c\alpha t)}{c\alpha}
\end{equation*}

\vsp

Set $g(x) = 0$ (thus $G(\alpha) = 0$) to find the solution for an initial displacement $f(x)$ with zero initial velocity:

\begin{align*}
    u(x,t) & = \frac{1}{2}\int_{-\infty}^\infty F(\alpha) e^{i\alpha(x+ct)} + F(\alpha) e^{i\alpha(x-ct)}~d\alpha \\
           & = \frac{1}{2} f(x + ct) + \frac{1}{2} f(x - ct)
\end{align*}

where Euler's formula yields

\begin{equation*}
    \cos(\alpha ct)e^{i\alpha x} = \frac{e^{i\alpha ct} + e^{-i\alpha ct}}{2}e^{i\alpha x}\text{.}
\end{equation*}

Now, set $f(x) = 0$ (thus $F(\alpha) = 0$) to find the solution for zero initial displacement with an initial velocity $g(x)$:

\begin{align*}
    u(x,t) & = \int_{-\infty}^\infty G(\alpha) \frac{\sin(\alpha ct)}{c\alpha} e^{i\alpha x}~d\alpha                                                \\
           & = \frac{1}{2c}\int_{-\infty}^\infty \frac{G(\alpha)}{i\alpha} e^{i\alpha(x+ct)} - \frac{G(\alpha)}{i\alpha} e^{i\alpha(x-ct)} ~d\alpha \\
           & = \frac{1}{2c} \int_{-\infty}^{x + ct}g(s)~ds - \frac{1}{2c} \int_{-\infty}^{x - ct}g(s)~ds                                            \\
           & = \frac{1}{2c} \int_{x - ct}^{x + ct}g(s)~ds
\end{align*}

where the Fourier transform's integration-division property yields the integrals in the spatial domain.

\vsp

Use the principle of superposition to combine the displacement and velocity solutions to find the general solution:

\begin{equation*}
    u(x,t) = \frac{1}{2}\left[f(x + ct) + f(x - ct)\right] + \frac{1}{2c} \int_{x - ct}^{x + ct}g(s)~ds
\end{equation*}

This is known as \textbf{d'Alembert's solution}.

\subsection{Convolution Theorem}

The convolution theorem states that multiplication in the spatial domain is equivalent to convolution in the frequency domain:

\begin{equation*}
    \mathcal{F}\left[f(x)g(x)\right] = F(\alpha) \ast G(\alpha)
\end{equation*}

where the convolution is defined as

\begin{equation*}
    f(x) \ast g(x) \equiv \int_{-\infty}^\infty f(\tau)g(x-\tau)~d\tau\text{.}
\end{equation*}

Note: this is assuming the definition of the Fourier transform presented in this reference. Since the placement of the constant $1/2\pi$ in the definition can vary, other versions may introduce an additional factor in the convolution theorem.

\vsp

The convolution thorem also goes the other way around: multiplication in the frequency domain is equivalent to convolution in the spatial domain:

\begin{equation*}
    \mathcal{F}\left[f(x) \ast g(x)\right] = 2\pi F(\alpha) G(\alpha)
\end{equation*}

Notice that the $2\pi$ factor appears in this second equation due to the definition of the Fourier transform used in this reference.